% 图契约：本题的源、边界和推导逐项对应；箭头符号按图注解释。
\documentclass[tikz,border=7pt]{standalone}
\usepackage{fontspec,xeCJK,amsmath,amssymb}
\IfFontExistsTF{Times New Roman}{\setmainfont{Times New Roman}}{\setmainfont{TeX Gyre Termes}}
\IfFontExistsTF{Heiti SC}{\setCJKmainfont[AutoFakeBold=1.5]{Heiti SC}}{\setCJKmainfont[AutoFakeBold=1.5]{FandolHei-Regular}}
\usetikzlibrary{arrows.meta,bending,calc,patterns,positioning}
\definecolor{ink}{HTML}{243745}
\definecolor{blue}{HTML}{275DAB}
\definecolor{green}{HTML}{177C63}
\definecolor{red}{HTML}{BD4935}
\definecolor{orange}{HTML}{AA6C16}
\definecolor{purple}{HTML}{715393}
\tikzset{every node/.style={font=\fontsize{9.3}{12.5}\selectfont,text=ink},
axis/.style={-{Stealth[length=4pt]},line width=.65pt,draw=ink!65},
vec/.style={-{Stealth[length=5pt]},line width=1pt,draw=blue},
short/.style={-{Stealth[length=3pt]},line width=.8pt},
guide/.style={dashed,draw=ink!40,line width=.5pt},
card/.style={rounded corners=2pt,draw=blue!65,fill=blue!4,align=center,text width=6.5cm,inner xsep=9pt,inner ysep=7pt},
title/.style={font=\fontsize{10.5}{14}\selectfont\bfseries,text width=15cm,align=center},
label/.style={font=\fontsize{8.4}{11}\selectfont,align=center},
note/.style={label,text width=6.2cm}}
\begin{document}
\begin{tikzpicture}[x=1cm,y=1cm]
\path[use as bounding box] (-8,-3.5) rectangle(8,3.5);
\node[title] at (0,3.15) {M2-146\quad 同号净电荷仍可能被极化吸引};
\begin{scope}[shift={(-4,0)}]
\filldraw[fill=ink!6,draw=ink!60](-.7,0)circle(1);\draw[guide](-1.8,0)--(2.4,0);\fill[ink](-.7,0)circle(1.4pt)node[above left]{$O$};\fill[red] (2.1,0) circle(2.3pt) node[below=3pt] {$q$};
\fill[blue] (-0.343,0) circle(2.3pt) node[below=3pt] {$q^{\prime}$};
\draw[<->,orange](-.7,1.45)--(2.1,1.45)node[midway,above]{$d$};\node[label]at(-1.5,-1.35){半径 $R$};\fill[green](-.7,0)circle(2.4pt);\draw[green](-1.2,.55)--(-.7,0);\node[label,green,above left]at(-1.2,.55){$q_0$};\node[note] at (0,-2.85) {临界量取决于距离，不能只按总电荷同号判断力向。};
\end{scope}
\node[card,draw=blue!65,fill=blue!4] at (3.9,1.8) {$\displaystyle Q_{0c}=\frac{qR^3(2d^2-R^2)}{d(d^2-R^2)^2}$};
\node[card,draw=green!65,fill=green!4] at (3.9,0) {$Q_0>Q_{0c}:\ \text{斥力}$\\$Q_0<Q_{0c}:\ \text{吸引}$};
\node[card,draw=orange!65,fill=orange!4] at (3.9,-1.8) {$d\to R^+:\quad Q_{0c}\to\infty$};
\end{tikzpicture}\end{document}